%!TEX root = ../thesis.tex

%jama %hakon -- joint discussion chapter

\chapter{Discussion}
\label{ch:discussion}

\section{Macro Model Assumptions and Limitations}
\label{sec:disc_macro_assumptions}

%hakon
The presented macro model is based on a set of simplifying assumptions that influence the interpretation of the results.
The model is not calibrated to represent a specific real-world outbreak, and most parameters are held constant.
As a result, the model should be interpreted as a theoretical framework rather than a predictive tool.

The population is modelled using normalised compartment values, representing proportions of the population rather than individuals.
This implies that infection levels can take small positive values without the epidemic dying out, which differs from stochastic models where the disease may naturally disappear at low case numbers.
In this model, this behaviour can be interpreted as representing undetected infections, but it also contributes to the emergence of repeated waves \citep{hethcote2000}.
Additionally, the model assumes a fixed population without demographic changes such as births, non-disease-related deaths, and age.
While this simplifies the dynamics, it limits the realism of long-term simulations.

Overall, these assumptions allow for a flexible modelling framework, but they also limit the accuracy of the results and should be considered when interpreting the behaviour of the model and the MPC controller.

\section{Macro Model Structure and Extensions}
\label{sec:disc_macro_structure}

%hakon
The model builds on the SIR and SEIR models that have been developed since 1927 \citep{kermack1927}.
Existing compartmental models with extended states and dynamics served as inspiration for the structure of the proposed model.
Some functions were simplified: age-based parameters were removed, and more complex structures such as susceptible-exposed-asymptomatic-symptomatic models were not included.
Vaccination was simplified to a constant rate; a more realistic approach would be time-dependent vaccination strategies such as batch vaccination at specific intervals.

To reduce the accumulation of vaccinated and immune individuals, waning immunity was introduced.
This is consistent with real-world diseases, COVID-19 being the primary motivation.
In the results section we saw that without waning immunity, there was no development of subsequent waves of infection.

The healthcare demand compartment was initially intended to represent hospitalisation strictly, with hospital bed availability as the control output for the MPC.
Because of the low reported number of hospital beds in Norway (approximately 5 beds per 100\,000 inhabitants \citep{kohler2021}), capacity was exceeded early in the simulation.
The $h$ parameter used was based on reported hospitalisation rates but may include more than just those who required hospital beds \citep{ariyaratne2025}.
The compartment was therefore interpreted as healthcare demand rather than strictly hospitalisation.

Both travel and commuting independently cause the cities to exhibit nearly identical compartment dynamics even when one starts with $I(0) = 0$.
This shows that the intercity mixing is strong.
This may resemble two cities in the same greater metropolitan area, such as Stavanger and Sandnes.
The travel matrix allows up to 2\% of the population to relocate every day, which is a high mobility rate for a typical urban setting but may be realistic for vacation hotspots.
Both the travel rate and the commuting matrix could have been modulated directly by the MPC in a more realistic approach; in the present formulation, the effect of the MPC on the transmission rate indirectly captures the impact of such interventions.

\section{MPC Behaviour and Control Performance}
\label{sec:disc_mpc}

%hakon
In Figure~\ref{fig:mpc_A_500}, the MPC is observed to maintain intervention for approximately 300~sample times in order to suppress the second and third wave shown in the uncontrolled 500-day scenario.
This results in a prolonged period of strong control action, which may appear excessive given that the measured output falls below the reference level during this period.

The behaviour is only weakly affected by the tuning parameter $w_u$ (manipulated variable weight), as similar responses were obtained for values of both 0.1 and 10.
This suggests that the observed behaviour is not solely due to controller tuning, but is also influenced by the underlying model and control formulation.

Several factors may contribute to this behaviour.
First, the choice of measured output may not fully capture the system dynamics relevant for control, and alternative outputs could potentially lead to different control actions.
Second, the epidemic model is highly nonlinear, while the MPC controller is based on a linearised approximation, which may reduce prediction accuracy outside the linearisation region.

Finally, the use of normalised compartment values implies that infection levels approach zero without reaching it exactly.
This results in a persistent low-level infection in the model, which may cause the controller to maintain intervention in order to prevent future growth.
This limitation suggests that alternative approaches, such as nonlinear or adaptive MPC \citep{grune2017}, could potentially improve performance by better capturing the underlying system dynamics.

\section{Choice of Controller Objective}
\label{sec:disc_objective}

%hakon
The measured output was defined as $H_{\mathrm{ratio}}$, representing the proportion of healthcare demand relative to available capacity.
This choice reflects the strategy used in Norway during the later stages of the COVID-19 pandemic, where avoiding excessive strain on the healthcare system was a primary objective \citep{stoustrup2023}.

However, the results indicate that $H_{\mathrm{ratio}}$ may not be an ideal control variable in the present model formulation.
The assumed healthcare capacity is relatively low, and with a hospitalisation rate of $h = 0.05$, the capacity is exceeded early in the simulation as the epidemic develops.
This leads to a control problem where the system operates close to or beyond capacity for much of the simulation.
Additionally, the hospitalisation rate used in the model may overestimate the demand for hospital resources, as reported values can include individuals admitted for reasons other than severe illness \citep{ariyaratne2025}.
These factors suggest that alternative or additional measured outputs, such as infection levels or combined indicators, could provide a more balanced basis for control.

\section{Interpretation of Cross-Scale Results}
\label{sec:disc_crossscale}

%jama
The simulation results support the main thesis of this project: that MPC-generated transmission rate signals, designed and validated at the macro (compartmental ODE) level, also effectively suppress epidemic dynamics when applied to a stochastic agent-based representation of the same population.
This cross-level consistency is not guaranteed a priori - the MPC controller has no knowledge of individual agent positions, travel decisions, or stochastic fluctuations - yet its output successfully guides the system toward lower infection and hospitalisation counts in the micro domain.

The explanation lies in the relationship between the population-level $\beta$ parameter and the individual-level infection probability.
When the MPC reduces $\beta(t)$, it reduces the per-contact infection probability for every agent in the corresponding city.
This blanket reduction is epidemiologically equivalent to enforcing social distancing (reducing contact frequency or contact intensity uniformly across the population), which is precisely the physical interpretation of reducing $\beta$ in the SIR framework.
The spatial proximity model in the ABM ensures that this reduction in per-contact probability translates to a reduction in the aggregate infection rate, recovering the macro-level dynamics in the mean-field limit.

This cross-level agreement is strengthened by the calibration procedure described in Chapter~\ref{ch:micro_results}.
Because the macro ODE and agent-based model operate at different levels of abstraction, agreement between them cannot be assumed even when nominal parameter values are matched.
The genetic-algorithm-based calibration aligns the macro parameters with the mean ABM hospital-load response, reducing discrepancy in the principal controlled variable of the MPC formulation.

\section{Value of the Multi-Scale Approach}
\label{sec:disc_multiscale}

%jama
The ABM reveals several features of epidemic dynamics that the macro model suppresses by construction \citep{polcz2025}.
Stochastic extinction events - where the epidemic dies out not because $R_t < 1$ but because the last infected agent recovers before transmitting - are possible in the ABM but impossible in the ODE model, which approaches the disease-free equilibrium only asymptotically.
Early-phase clustering effects, where infection spreads rapidly in a local neighbourhood before dispersing city-wide, produce heterogeneous spatial patterns not visible in aggregate trajectories.
Inter-city transmission via commuting creates bursty coupling events that the contact matrix $M$ only approximates in expectation.

These features matter for public health practice.
A deterministic macro model may predict hospital load smoothly approaching a threshold, while the stochastic micro model may show occasional realisations where hospital load spikes suddenly due to a single superspreading event or a large commuting cluster.
Risk-aware policy design should account for this tail behaviour.

\section{Model Applicability and Future Work}
\label{sec:future}

%jama %hakon
The framework developed here is modular and suitable for a range of extensions.
The following directions are identified as most promising:

\paragraph{Bidirectional coupling.}
Implement online feedback from the ABM to the MPC controller, using filtered aggregate statistics from the micro model as the observed state.
This would require either fast ABM implementations or surrogate models approximating ABM behaviour \citep{grune2017,polcz2025}.
Physics-informed neural network approaches \citep{zhong2025} could serve as efficient surrogate models for this purpose.

\paragraph{Population scale.}
Extend the ABM to larger populations using efficient data structures (vectorised MATLAB operations, GPU parallelism) or agent approximation methods \citep{kerr2021}.
This would reduce stochastic variance and produce more reliable mean-field convergence.

\paragraph{Network-based contacts.}
Replace the spatial proximity model with an explicit contact network (household, workplace, school layers) informed by empirical contact matrices \citep{zhang2021} to better represent heterogeneous transmission \citep{grundel2022}.

\paragraph{Robust MPC under uncertainty.}
Incorporate uncertainty in the travel matrix $M$ and in the contact rates following the framework of \citet{albi2021}, producing intervention policies that are robust to misspecified social mixing patterns.

\paragraph{Multi-city extension.}
Generalise the two-city framework to a network of multiple cities with heterogeneous sizes and connection strengths, enabling study of epidemic spread across realistic metropolitan or national geographies \citep{carli2020}.

\paragraph{Uncertainty quantification.}
Run multiple stochastic realisations of the ABM under the same $\beta(t)$ input to quantify the distribution of outcomes and evaluate the robustness of the MPC policy to stochastic variability \citep{peak2017}.

\paragraph{Data-driven parameter calibration.}
Apply genetic-algorithm-based parameter estimation \citep{rajaei2021} to fit the model to Norwegian COVID-19 incidence data, producing a calibrated simulation environment suitable for retrospective policy evaluation and prospective planning \citep{ariyaratne2025,dong2022}.

\paragraph{Nonlinear or adaptive MPC.}
Develop a nonlinear MPC formulation or an adaptive variant that updates its internal model based on observed epidemic state, reducing the performance loss caused by the linearisation approximation \citep{morato2020,liu2021}.
